% Numerical Setup for SN_A / SN_B Double-Explosion Modelling
% SN2024pba Project
%
% Compile: module load texlive/2021
%          pdflatex numerical_setup && bibtex numerical_setup && pdflatex numerical_setup && pdflatex numerical_setup

\documentclass[11pt, a4paper]{article}

\usepackage[round, sort, compress]{natbib}
\usepackage{amsmath, amssymb}
\usepackage{graphicx}
\usepackage{booktabs}
\usepackage{hyperref}
\usepackage[T1]{fontenc}
\usepackage{microtype}

% Shortcuts
\newcommand{\Msun}{\ensuremath{M_\odot}}
\newcommand{\foe}{\ensuremath{10^{51}\;\mathrm{erg}}}
\newcommand{\kms}{\ensuremath{\mathrm{km\;s^{-1}}}}
\newcommand{\cgs}{\ensuremath{\mathrm{g\;cm^{-3}}}}
\newcommand{\kappaes}{\ensuremath{\kappa_{\mathrm{es}}}}
\newcommand{\taumin}{\ensuremath{\tau_{\mathrm{min}}}}
\newcommand{\taumax}{\ensuremath{\tau_{\mathrm{max}}}}
\newcommand{\Ncnd}{\ensuremath{N_{\mathrm{cnd}}}}
\newcommand{\FitTau}{\ensuremath{\mathtt{FitTau}}}
\newcommand{\NCNMIN}{\ensuremath{N_{\mathrm{CNMIN}}}}
\newcommand{\EKO}{\ensuremath{E_{\mathrm{KO}}}}
\newcommand{\Ni}{\ensuremath{^{56}\mathrm{Ni}}}
\newcommand{\Co}{\ensuremath{^{56}\mathrm{Co}}}
\newcommand{\Fe}{\ensuremath{^{56}\mathrm{Fe}}}
\newcommand{\sigmav}{\ensuremath{\sigma}}
\newcommand{\vwind}{\ensuremath{v_{\mathrm{wind}}}}
\newcommand{\mdot}{\ensuremath{\dot{M}}}
\newcommand{\Rstar}{\ensuremath{R_\star}}
\newcommand{\Lstar}{\ensuremath{L_\star}}
\newcommand{\tdays}{\ensuremath{t_{\mathrm{days}}}}
\newcommand{\tSBO}{\ensuremath{t_{\mathrm{SBO}}}}

\begin{document}

\title{Numerical Setup for the SN\_A / SN\_B Double-Explosion Modelling of SN\,2024pba}

\author{S. Brennan\\
\small Max Planck Institute for Astrophysics, Karl-Schwarzschild-Str.\ 1, 85741 Garching, Germany}

\section{Numerical Setup}
\label{sec:numerical_setup}

This section describes the numerical methods and model configuration used to
simulate the double-explosion scenario consisting of a Type~IIb supernova
(SN\_A) and a subsequent H/He shell ejection into the SN\_A ejecta (SN\_B).
The modelling pipeline couples the multi-group radiation hydrodynamics code
STELLA \citep{Blinnikov1998, Blinnikov2000, Blinnikov2006} with the non-LTE
radiative transfer code CMFGEN \citep{Hillier1998, HillierDessart2012,
Dessart2015} for spectral synthesis at selected epochs.

\subsection{Overview of the Modelling Pipeline}
\label{subsec:pipeline_overview}

The simulation proceeds in two stages.  In the first stage (SN\_A), a
stripped-envelope progenitor is exploded as a Type~IIb supernova and evolved
with STELLA from shock breakout (SBO) through the radioactive-decay-powered
light-curve phase.  A snapshot at $\sim$30\,d post-SBO is extracted and
remapped onto a CMFGEN depth grid for non-LTE spectral synthesis.  The SN\_A
ejecta is then homologously evolved to late times ($\sim$400\,d) to serve as
the circumstellar medium (CSM) for the second stage.

In the second stage (SN\_B), a H/He shell is placed at the centre of the
homologously expanded SN\_A ejecta and launched at 5000\,$\kms$ into the
surrounding CSM.  STELLA evolves the shell--CSM interaction, and a snapshot
at $\sim$20--40\,d after shell launch is extracted for CMFGEN spectral
synthesis.

The entire pipeline is orchestrated by a single workflow manager
(\texttt{run\_model\_grid.py}) that handles progenitor conversion, STELLA
job submission, snapshot extraction, CMFGEN grid generation, and
dependency-chained spectral synthesis jobs on a high-performance computing
(HPC) cluster.

\subsection{SN\_A: Type~IIb Explosion of a Stripped Helium Star}
\label{subsec:sn_a}

\subsubsection{Progenitor Model and Pre-SN Conversion}
\label{subsubsec:progenitor}

The SN\_A progenitor is based on the solar-metallicity $12\,M_\odot$ model
\texttt{s12.0} from \citet{WoosleyHeger2007}, evolved through core collapse
with the KEPLER stellar evolution code.  The hydrogen envelope is removed to
mimic binary stripping, leaving a $\sim$3.0\,$\Msun$ helium core.  This
configuration matches model~12C of \citet{Ergon2022}, which was constructed to
reproduce the Type~IIb SN\,2011dh \citep{Ergon2015}.

The KEPLER pre-SN model is converted to STELLA's Lagrangian hydrodynamics
format by the script \texttt{convert\_kepler\_presn\_to\_stella.py}, which
performs the following operations:
\begin{enumerate}
    \item \textbf{Mass cut excision:} The innermost 1.4\,$\Msun$ (the compact
    remnant) is removed, leaving $\sim$1.6\,$\Msun$ of ejecta.
    \item \textbf{Composition remapping:} KEPLER isotopes are mapped onto
    STELLA's 15-element nuclear network (H, He, C, N, O, Ne, Na, Mg, Al, Si,
    S, Ar, Ca, Fe, Ni).  The iron-group isotopes ($^{56}$Fe, $^{56}$Ni,
    $^{44}$Ti, $^{40}$Ca, $^{36}$Ar, $^{32}$S) are combined into a single
    Fe-peak tracer.
    \item \textbf{H surface layer addition:} A thin hydrogen envelope of
    0.03--0.10\,$\Msun$ with H mass fraction $X = 0.7$ is added to the
    outermost zones, reproducing the residual H layer characteristic of
    Type~IIb progenitors \citep{Ergon2022}.
    \item \textbf{CSM wind attachment:} A low-density $1/r^2$ circumstellar
    wind is appended beyond the stellar surface using a sigmoid-blended
    transition that ensures $C^\infty$ continuity in both the density and
    all its derivatives.  The wind has a mass-loss rate
    $\mdot = 10^{-6}\,\Msun\,\mathrm{yr^{-1}}$, a terminal velocity
    $\vwind = 100\,\kms$, and a duration of 320\,yr, yielding an outer radius
    of $\sim$$10^{17}$\,cm and a total CSM mass of $\sim$$3 \times
    10^{-4}\,\Msun$.  The CSM is primarily required for numerical stability
    --- it prevents the STELLA radiation solver from stalling when the
    photosphere exits the compact ejecta at $\sim$10\,d --- while contributing
    negligible mass and optical depth to the light curve.
    \item \textbf{Density gradient smoothing:} Steep density gradients
    ($|d\log\rho / d\log r| > 5$) in the pre-SN profile are smoothed with a
    Gaussian kernel ($\sigma = 3$ zones) to prevent GEAR solver stalls at
    the initial condition.
\end{enumerate}

The resulting profile is remapped onto 500 Lagrangian zones (450 ejecta
zones + 50 wind zones) and written in STELLA's \texttt{.hyd} (hydrodynamics)
and \texttt{.abn} (abundance) binary formats.

\subsubsection{STELLA Radiation Hydrodynamics}
\label{subsubsec:stella_sn_a}

STELLA \citep{Blinnikov1998, Blinnikov2000, Blinnikov2006} solves the
one-dimensional Lagrangian radiation hydrodynamics equations with
multi-group grey radiative transfer in the homologous comoving frame.  The
code employs a variable Eddington factor closure and treats radiation
transport in two regimes: a diffusion approximation for the optically
thick interior ($\tau > \FitTau$) and a ray-by-ray (characteristic) method
for the optically thin outer layers ($\tau < \FitTau$).  The boundary
between these regimes, $\Ncnd$, is determined at each timestep by searching
for the depth where the optical depth at the minimum-opacity frequency
equals $\FitTau$.

\paragraph{Explosion mechanism.}
The explosion is driven by a thermal bomb deposited in the inner
0.3\,$\Msun$ of the ejecta over a duration of 10\,s, with a total energy
input of $0.9 \times 10^{51}$\,erg.  The shorter bomb duration (10\,s
versus the 100\,s used in earlier runs) increases the kinetic-energy
conversion efficiency from $\sim$61\% to $\sim$80\%, yielding a final
kinetic energy of $\sim$0.7\,foe consistent with the Ergon et~al.\ (2022)
value for SN\,2011dh.  The parameter $\EKO = 0$ is set so that STELLA
preserves the zero initial velocities from the \texttt{.hyd} file; the
thermal bomb generates the explosion and establishes the velocity profile
during the first $\sim$0.1\,d of evolution.

\paragraph{Radioactive heating.}
A \Ni\ mass of 0.075\,$\Msun$ \citep{Ergon2022} is distributed uniformly
over the inner 0.3\,$\Msun$ of the ejecta (the inner 40\% of zones).  The
\Ni\ is read from a dedicated \texttt{.xni} file and injected into the
composition via the \texttt{eve2} pre-processor.  The radioactive decay
chain $\Ni \to \Co \to \Fe$ is followed self-consistently by STELLA's
nuclear network, with gamma-ray deposition treated using the
\citet{Swartz1995} opacity formalism.

\paragraph{Numerical parameters.}
The key STELLA control parameters for the SN\_A model are summarised in
Table~\ref{tab:stella_params}.  The GEAR integrator \citep{Gear1971,
ByrneHindmarsh1975} is used with a relative tolerance of
$\epsilon = 3 \times 10^{-3}$, a minimum timestep of $10^{-14}$\,s, and a
maximum timestep of $\sim$5.8\,d.  The artificial viscosity parameters
($A_Q = 50$, $B_Q = 1$) are set high to damp the thermal-bomb shock at
early times ($t < 1$\,d); lower values ($A_Q = 20$) caused matrix solver
failures at $t \sim 7 \times 10^{-5}$\,d.

The $\FitTau$ parameter is set to 5.0, following the STELLA reference
model \texttt{X\_scat2.dat}.  This value is appropriate for SN ejecta with
strongly frequency-dependent opacity, where the optical depth at the
minimum-opacity frequency is $\sim$100$\times$ smaller than the grey
optical depth.  Setting $\FitTau = 15$ (as in the homogeneous test model
\texttt{m030307mhomo.dat}) causes the $\Ncnd$ boundary to collapse
prematurely when the minimum-frequency $\tau$ drops below 15, leading to a
catastrophic jump in the diffusion--ray-by-ray boundary and a failure of
the Y12M sparse matrix solver.

To prevent $\Ncnd$ collapse at late times ($t > 10$\,d), where the
minimum-frequency optical depth inevitably drops below any reasonable
$\FitTau$ due to geometric expansion, the minimum number of diffusion
zones is set to $\NCNMIN = N_{\mathrm{zon}} / 10 = 50$.  Additionally, the
per-timestep change in $\Ncnd$ is limited to $\Delta N_{\mathrm{cnd}} = 20$
zones, converting catastrophic single-step jumps (e.g.\ 279$\to$50) into a
gradual descent over $\sim$12 timesteps that the matrix solver can
tolerate.

\begin{table}[ht]
\centering
\caption{STELLA numerical parameters for the SN\_A and SN\_B models.}
\label{tab:stella_params}
\begin{tabular}{lcc}
\toprule
Parameter & SN\_A & SN\_B \\
\midrule
Number of zones ($N_{\mathrm{zon}}$) & 500 & 500 \\
Target duration (d) & 100 & 45 \\
Explosion energy ($10^{51}$\,erg) & 0.9 (thermal bomb) & 0 (kinetic) \\
$\EKO$ & 0 & 0 \\
$A_Q$ (artificial viscosity) & 50 & 10 \\
$B_Q$ (pseudo-viscosity) & 1 & 1 \\
$\FitTau$ & 5.0 & 5.0 \\
$\NCNMIN$ & $N_{\mathrm{zon}}/10$ & $N_{\mathrm{zon}}/10$ \\
$\Delta N_{\mathrm{cnd}}$ limit & 20 & 20 \\
GEAR method & 2 & 3 \\
$\epsilon$ (integrator tolerance) & $3 \times 10^{-3}$ & $3 \times 10^{-3}$ \\
$h_{\min}$ (s) & $10^{-14}$ & $10^{-14}$ \\
$h_{\max}$ (s) & $5 \times 10^{5}$ & 3600 \\
$T_{\mathrm{floor}}$ (K) & 300 & 300 \\
$T_{\mathrm{rad,\,floor}}$ (K) & 100 & 100 \\
\Ni\ mass ($\Msun$) & 0.075 & 0 \\
Scattering & Enabled & Enabled \\
\bottomrule
\end{tabular}
\end{table}

\subsubsection{Homologous Evolution to Late Times}
\label{subsubsec:homologous}

At late times ($t \gtrsim 60$\,d), the SN\_A ejecta is in free expansion
($v = r/t$), and the evolution is accurately described by homologous
scaling.  The script \texttt{homologous\_evolution.py} scales the final
STELLA snapshot to a target age $t_{\mathrm{target}}$ using:
\begin{align}
    r(t) &= r(t_0) \, \frac{t}{t_0}, \\
    v(t) &= v(t_0) \quad (\text{constant}), \\
    \rho(t) &= \rho(t_0) \left(\frac{t_0}{t}\right)^3, \\
    T(t) &= T(t_0) \, \frac{t_0}{t} \quad (\text{adiabatic}),
\end{align}
where $t_0$ is the age of the last STELLA epoch.  These relations are valid
when the expansion is highly supersonic, there is no ongoing energy
injection (the \Ni\ decay chain is negligible at $t \gtrsim 200$\,d), and
there is no significant CSM interaction.  The homologously evolved profile
at 400\,d serves as the CSM for the SN\_B model.

\subsection{SN\_B: H/He Shell Interaction with the SN\_A Ejecta}
\label{subsec:sn_b}

\subsubsection{Initial Model Construction}
\label{subsubsec:sn_b_initial}

The SN\_B initial model is constructed by the script
\texttt{prepare\_sn\_b.py} (or \texttt{prepare\_x\_model.py} for the
three-component variant with a connecting wind).  The model consists of
two (or three) components, assembled from inside out:

\begin{enumerate}
    \item \textbf{H/He shell (inner):} A uniform-density sphere of
    1.0\,$\Msun$ with composition $X = 0.7$, $Y = 0.3$ and a homologous
    velocity profile reaching 5000\,$\kms$ at its outer edge.  The shell
    represents material ejected by a secondary explosion or a
    pulsational mass-loss event.  The density structure follows a
    power-law $\rho \propto r^{-n_\rho}$ with $n_\rho = 8$, consistent
    with the reverse-shock shell models of \citet{Dessart2016} for
    SN\,1994W.  The temperature profile follows $\rho \propto r^{-n_T}$
    with $n_T = 0.4$, with a maximum core temperature of 65\,000\,K.

    \item \textbf{Connecting wind (middle, optional):} A $1/r^2$ wind
    with $\mdot = 0.1\,\Msun\,\mathrm{yr^{-1}}$ and $\vwind = 100\,\kms$
    bridges the shell and the old ejecta, preventing a hard density
    discontinuity that would stall the GEAR solver.  This component is
    included in the three-component X~model variant.

    \item \textbf{SN\_A ejecta (outer):} The 400-day homologously evolved
    SN\_A profile (\S\ref{subsubsec:homologous}), which now acts as the
    CSM.  The ejecta retains its homologous velocity structure and
    He-rich, metal-enhanced composition from the Type~IIb explosion.
\end{enumerate}

The components are combined into a single 500-zone STELLA
\texttt{.hyd} file.  The parameter $\EKO = 0$ is essential here: it
preserves both the shell velocity (5000\,$\kms$) and the CSM homologous
velocity profile.  Setting $\EKO > 0$ would overwrite all velocities with
a triangular (linear-in-mass) profile, destroying the physically motivated
shell--CSM velocity contrast and causing the solver to stall.

\subsubsection{STELLA Evolution of the Interaction}
\label{subsubsec:stella_sn_b}

STELLA evolves the SN\_B model for 45\,d from the shell launch time.
The H/He shell, moving at 5000\,$\kms$, drives a strong forward shock
into the slower CSM ($v \sim 500$--$6000\,\kms$), while a reverse shock
propagates back into the shell.  The shock interaction converts kinetic
energy into radiation, producing the interaction-powered luminosity that
characterises Type~IIn supernovae \citep{Dessart2016, Chugai2001}.

The artificial viscosity is set to $A_Q = 10$ (higher than the SN\_A
value of 2--4 for a standard thermal-bomb explosion) to stabilise the
GEAR solver across the strong velocity reversal at the shell--CSM
interface.  The GEAR method is set to 3 (GEAR+BGH) for improved stability
in the stiff interaction system.  No \Ni\ is included in the SN\_B model;
the luminosity is entirely interaction-powered.  The maximum timestep is
limited to 3600\,s (1\,h) to resolve the early shock interaction phase.

\subsection{CMFGEN Non-LTE Spectral Synthesis}
\label{subsec:cmfgen}

\subsubsection{STELLA-to-CMFGEN Conversion}
\label{subsubsec:conversion}

At selected epochs (30\,d post-SBO for SN\_A; 20--40\,d post-launch for
SN\_B), a snapshot is extracted from the STELLA \texttt{.swd} output file
and converted to CMFGEN's \texttt{SN\_HYDRO\_DATA} format by the script
\texttt{convert\_stella\_swd\_to\_cmfgen.py}.  The conversion performs
the following steps:

\begin{enumerate}
    \item \textbf{Epoch selection:} The STELLA epoch closest to the
    requested snapshot day (within a tolerance of $\pm$2\,d) is selected.
    If no STELLA epoch is available (e.g.\ for late-time snapshots beyond
    the STELLA target), the homologous evolution procedure
    (\S\ref{subsubsec:homologous}) is applied.

    \item \textbf{Optical-depth filtering:} Zones are filtered to retain
    only the region $\taumin \leq \tau \leq \taumax$ (typically
    $0.01 \leq \tau \leq 30$), where $\tau$ is computed from the
    STELLA NLTE opacity.  This removes both the optically thin vacuum
    exterior and the optically thick interior where CMFGEN's
    Sobolev approximation breaks down.

    \item \textbf{Adaptive grid remapping:} The filtered STELLA zones are
    remapped onto a CMFGEN depth grid using an adaptive algorithm that
    concentrates zones in regions of steep gradients (ionisation fronts,
    density jumps, velocity reversals).  The target grid size is
    $N_D \sim 150$--200 points, with constraints on maximum cell-to-cell
    jumps in radius ($\Delta r / r \leq 0.05$), temperature
    ($\Delta \log T \leq 0.10$), density ($\Delta \log \rho \leq 0.15$),
    and velocity ($\Delta v \leq 0.30\,v$).  The velocity--density
    structure parameter $\sigmav = d\ln v / d\ln r - 1$ is monitored, and
    additional zones are inserted where $|\sigmav|$ or $|d\sigmav / dr|$
    exceeds threshold values, ensuring adequate resolution of the
    Sobolev line optical depth.

    \item \textbf{Composition transfer:} The STELLA 15-element composition
    is mapped onto CMFGEN's species list.  The \Ni\ distribution is read
    from the STELLA \texttt{.xni} file (Fortran unformatted REAL*8) and
    included in the \texttt{SN\_HYDRO\_DATA} abundance table so that
    CMFGEN can compute radioactive decay heating self-consistently.

    \item \textbf{Velocity smoothing:} Small numerical oscillations in
    the STELLA velocity profile ($\pm 0.3\,\kms$) are smoothed with a
    Savitzky--Golay filter (window $= 5$ zones, $\Delta v = 50\,\kms$)
    to prevent the structure parameter $\sigmav$ from reaching unphysical
    values ($\sigmav < -19$) that destabilise the CMFGEN Sobolev
    integrals.
\end{enumerate}

The output \texttt{SN\_HYDRO\_DATA} file is written in CMFGEN's
outer-to-inner (decreasing radius) depth convention, with radius in units
of $10^{10}$\,cm, velocity in $10^5\,\mathrm{cm\,s^{-1}}$, temperature in
$10^4$\,K, and density in $\mathrm{g\,cm^{-3}}$.

\subsubsection{CMFGEN Configuration and Convergence Strategy}
\label{subsubsec:cmfgen_config}

CMFGEN \citep{Hillier1998, HillierDessart2012, Dessart2015} computes
non-LTE level populations and emergent spectra for a spherically symmetric
supernova atmosphere.  The code solves the radiative transfer equation in
the Sobolev approximation with continuum opacities treated in the
co-moving frame, and iterates the statistical equilibrium equations
coupled to the radiative transfer via an accelerated Lambda iteration
(ALI) scheme \citep{RybickiHummer1991}.

\paragraph{Velocity law and grid.}
For both SN\_A and SN\_B, the velocity law is set to
\texttt{VEL\_LAW\,=\,12}, which reads the actual (non-homologous)
velocity profile from the \texttt{SN\_HYDRO\_DATA} input rather than
fitting a homologous $v \propto r$ relation.  This is essential for the
SN\_B model, where the shell--CSM velocity reversal produces
$\sigmav < -1$ regions that a homologous fit cannot represent.  The
\texttt{VEL\_LAW\,=\,12} routine includes a protective clamp that sets
$\sigmav = -0.999$ when $-1.2 < \sigmav < -1.0$, preventing the Sobolev
line optical depth denominator $|1 + \sigmav \mu^2|$ from approaching
zero.  Empirical tests confirm that removing this clamp causes
catastrophic divergence of the Ng acceleration extrapolation, even though
the \citet{Dessart2015} Sobolev guard
$\max(|1 + \sigmav \mu^2|, -0.2(1+\sigmav), 0.1)$ is algebraically
well-defined for all finite $\sigmav$.

\paragraph{Solver configuration.}
The depth-coupling solver is set to \texttt{TRIDIAG}
(\texttt{NUM\_BNDS\,=\,3}), which allows inter-depth error propagation
via a tridiagonal band matrix.  This is required for models with
non-monotonic velocity ($\sigmav < -1$), where the diagonal-only
(\texttt{DIAG}) solver cannot capture the inter-depth coupling introduced
by the velocity reversal.  The formal solution uses
\texttt{N\_TYPE\,=\,N\_ON\_J} with \texttt{USE\_LAM\_ES\,=\,T}, which
sets the mean intensity $J$ to the formal-solution value $J_{\mathrm{TC}}$
and bypasses the moment equations, improving stability for models with
large velocity gradients.

\paragraph{Iteration strategy.}
The convergence strategy employs alternating LAMBDA and Boltzmann
Acceleration (BA) iterations, controlled by the following VADAT
parameters:
\begin{itemize}
    \item \texttt{DO\_NG\,=\,F}: Ng acceleration \citep{Ng1974} is
    disabled.  While Ng acceleration can speed convergence for well-behaved
    models, it destabilises trace species (e.g.\ Fe~\textsc{iv},
    Fe~\textsc{v}) at ionisation fronts by extrapolating across
    oscillatory population changes, leading to divergence of up to
    $10^{26}\%$ in the maximum correction.
    \item \texttt{DO\_AV\,=\,T}: Population averaging
    (\texttt{AVE\_FLIPS}) is enabled to damp oscillating populations
    between LAMBDA iterations.  A source-code patch was applied to allow
    \texttt{AVE\_FLIPS} to fire during continuous LAMBDA iterations
    (\texttt{DO\_LAM\_IT\,=\,T}), which was previously blocked by a
    condition that required \texttt{MAIN\_COUNTER $>$ LAST\_LAMBDA +
    NUM\_OSC\_AV + 2} --- a condition that is never satisfied when every
    iteration is a LAMBDA iteration.
    \item \texttt{MAX\_LAM\,=\,$10^5$}: The maximum fractional population
    change per LAMBDA iteration is capped at $10^5$, with
    \texttt{LAM\_SCALE\_OPT\,=\,LIMIT} applying a proportional limiter to
    trace species.
    \item \texttt{MAX\_LIN\,=\,3}: The maximum number of linearisation
    steps per BA iteration is limited to 3 (reduced from the manual
    default of 10) for cold-start stability, as the BA solve can produce
    extreme corrections when the initial LTE populations are far from the
    converged solution.
    \item \texttt{BA\_CHK\_FAC\,=\,0.05}: The BA matrix is held fixed
    when the maximum population correction falls below 5\%, avoiding
    unnecessary matrix recomputations.
    \item \texttt{METHOD\,=\,ZERO}: The trapezoidal optical-depth
    integration method is used, as recommended by the CMFGEN
    documentation for models with convergence difficulties.
    \item \texttt{FIX\_T\,=\,T}: The temperature is fixed to the
    hydrodynamic (STELLA) value at all depths throughout the iteration,
    preventing CMFGEN from solving for temperature and using the
    physically consistent hydrodynamic temperature instead.
\end{itemize}

\paragraph{Inner boundary condition.}
The inner boundary condition is set to \texttt{DIFFUSION}
(\texttt{DIF\,=\,T}, \texttt{IB\_METH\,=\,DIFFUSION}) when the inner
optical depth $\tau_{\mathrm{inner}} \geq 50$, providing a physically
motivated radiation field seed ($\Lstar$) from the STELLA luminosity.
When $\tau_{\mathrm{inner}} < 50$, the \texttt{DIFFUSION} boundary is
still used (with a warning), as a non-zero $\Lstar$ is preferable to the
\texttt{ZERO\_FLUX} alternative, which provides no radiation seeding and
produces $L = 0$ at iteration~1.

\paragraph{Species selection.}
The atomic model includes H~\textsc{i}, He~\textsc{i}, O~\textsc{i},
Si~\textsc{i--ii}, Ca~\textsc{i--ii}, and Fe~\textsc{i--iii} as the
primary spectral-line species.  High-ionisation iron species
(Fe~\textsc{iv}, Fe~\textsc{v}) are removed for the 30\,d SN\_A snapshot,
as their ionisation potentials ($\chi = 54.8$ and 77\,eV) are far above
the maximum ejecta temperature ($T_{\max} \approx 7700$\,K at 30\,d),
making their NLTE populations $\sim$$10^{-26}$ (set by photoionisation
from the photospheric radiation field) and introducing $10^{16}\%$
trace-species noise into the convergence.  Similarly, C~\textsc{ii} is
reduced to ground-state-only (\texttt{ISF\,=\,1,1,1}) for models with
near-zero carbon abundance ($X_C \sim 10^{-15}$), as 8 NLTE levels with
zero abundance drove $10^{36}$ fractional population cascades.

\paragraph{Spectral synthesis.}
Once the level populations have converged, CMFGEN computes the emergent
spectrum via a formal solution of the radiative transfer equation on a
fine frequency grid spanning 900--28\,000\,\AA\ (covering the optical
through mid-infrared for JWST MIRI comparison).  The spectrum includes
line blanketing from all included species, continuum opacities
(bound-free, free-free, electron scattering), and the diffuse emission
from the SN ejecta.  An observed-frame spectrum is produced by applying
the cosmological redshift and distance modulus appropriate to the target
SN.

\subsection{Computational Infrastructure}
\label{subsec:infrastructure}

All simulations are performed on the MPCDF Viper HPC cluster using SLURM
job scheduling.  STELLA runs on a single node with 16\,GB memory and a
24\,h wall-time limit.  CMFGEN runs use the same resources, with
dependency-chained INIT$\to$STABLE$\to$FINAL or single-FINAL
configurations.  The workflow manager \texttt{run\_model\_grid.py}
handles automatic job cancellation of previous runs, directory cleanup,
VADAT/MODEL\_SPEC regeneration from YAML configuration, and consistent
validation of solver parameters across all model components.

\section*{Acknowledgments}
This research made use of the MPCDF Viper cluster at the Max Planck
Computing and Data Facility.  We acknowledge the use of the KEPLER,
STELLA, and CMFGEN codes, and thank S.\ Woosley, A.\ Heger,
S.\ Blinnikov, and D.\ Hillier for making these codes available to the
community.

\paragraph{Software}
STELLA \citep{Blinnikov1998, Blinnikov2000, Blinnikov2006},
CMFGEN \citep{Hillier1998, HillierDessart2012, Dessart2015},
KEPLER \citep{WoosleyHeger2007},
NumPy \citep{Harris2020},
Matplotlib \citep{Hunter2007},
SciPy \citep{Virtanen2020}.

\bibliographystyle{aasjournal}
\bibliography{references}

\end{document}
